En este trabajo se abordó el reto de parametrizar el modelo de sinapsis química gradual de \cite{Golowasch1999} para su aplicación en circuitos biohíbridos. Hallar parámetros que reproduzcan dinámicas específicas (como el comportamiento en fase o antifase) es complejo debido al elevado número de los mismos, a su acoplamiento, y al fenómeno de \textit{parameter sloppiness}. A esta complejidad se suman las estrictas restricciones temporales de la experimentación con neuronas vivas, donde el proceso de ajuste no solo debe realizarse con rapidez para garantizar la viabilidad del tejido, sino que, en caso de ser automático, sus procesos subyacentes no pueden comprometer los plazos temporales críticos exigidos por el \ac{rtos}.

Para superar estas limitaciones, se desarrollaron dos arquitecturas basadas en \ac{bo}: una herramienta \textit{offline} para validaciones rápidas y un módulo \textit{online} integrado en \ac{rtxi}. El uso de un \ac{gp} con \textit{kernel} \ac{se} provisto de \ac{ard} se justificó por su eficiencia muestral y capacidad para ponderar la relevancia de cada parámetro, mitigando autónomamente el solapamiento funcional. El espacio de búsqueda se definió cuidadosamente, adaptando los rangos a cada experimento e incorporando reparametrizaciones (uso de $R$) y escalas logarítmicas para contrarrestar el \textit{parameter sloppiness}. A nivel técnico, la implementación \textit{online} destacó por su arquitectura de sincronización \textit{lock-free} entre el hilo de la \ac{bo} y el hilo de tiempo real, garantizando el cálculo ininterrumpido de la sinapsis.

La validación técnica corroboró la eficacia del algoritmo, mostrando una convergencia rápida hacia regiones óptimas del espacio de parámetros. La función de evaluación diseñada, que separa métricas de rango y de forma (estas últimas basadas en la correlación de Pearson) orientó eficazmente la búsqueda, reproduciendo con precisión las componentes de baja (onda lenta) y alta frecuencia (\textit{spikes}) del potencial presináptico. Las simulaciones confirmaron que, con solo 400 evaluaciones y la adquisición guiada por \ac{ei}, el algoritmo estabilizó los resultados de forma consistente, superando al muestreo aleatorio \ac{lhs} inicial.

Empíricamente, las pruebas funcionales ratificaron la validez de las implementaciones generadas. Tanto en simulaciones \textit{offline} unidireccionales (empleando registros de neuronas vivas y modelos de \ac{hr}, resueltas en apenas medio minuto al procesar datos estáticos), como en evaluaciones \textit{online} bidireccionales (con dos modelos de \ac{hr}), el sistema indujo con solvencia interacciones de fase y antifase mediante unas corrientes con una forma realista. Optimizando simultáneamente ambas contribuciones ($I_{\mathrm{fast}}$ e $I_{\mathrm{slow}}$) en las simulaciones \textit{offline}, y parametrizando $I_{\mathrm{fast}}$ en una dirección e $I_{\mathrm{slow}}$ en la opuesta durante las pruebas \textit{online}, se obtuvieron los acoplamientos deseados en menos de 16 minutos para el entorno \textit{online}, un tiempo plenamente compatible con la viabilidad de los experimentos con neuronas vivas.

Esta herramienta augura un impacto muy favorable en futuras investigaciones que requieran acoplamientos complejos, tanto en contextos biohíbridos como computacionales. Al automatizar el tedioso ajuste manual, el sistema reduce la carga preparatoria de los experimentos y ofrece gran flexibilidad para estudiar topologías donde las componentes frecuenciales se modulen por separado. Además, el código y la metodología están plenamente integrados en la librería \ac{rthybrid} sobre \ac{rtxi} 2.3, asegurando su disponibilidad inmediata para las configuraciones experimentales del \ac{gnb} de la \ac{uam}.

Pese a cumplir los objetivos planteados, el desarrollo presentó ciertas limitaciones. El módulo \textit{online} es específico para \ac{rtxi} 2.3 (para asegurar su compatibilidad inminente con el laboratorio del \ac{gnb}), lo que requerirá adaptaciones para funcionar en versiones más recientes. Por otro lado, el ajuste paramétrico mostró ligeras imprecisiones al reproducir los extremos del rango de corriente exactos. Estas desviaciones probablemente se debieron a rangos objetivo imposibles de alcanzar exactamente en general o sin sacrificar la puntuación de la forma de onda. No obstante, este pequeño margen de error no impidió el correcto funcionamiento del algoritmo ni la obtención de los invariantes dinámicos perseguidos.

Como trabajo futuro, el paso inmediato es validar el módulo \textit{online} con neuronas biológicas del ganglio estomatogástrico en el laboratorio del \ac{gnb}, comprobando su adaptación a la variabilidad biológica real y la idoneidad de los tiempos de optimización. También resultaría interesante optimizar simultáneamente ambas componentes ($I_{\mathrm{fast}}$ e $I_{\mathrm{slow}}$) en las dos direcciones de un circuito bidireccional, evaluando cómo el aumento de dimensionalidad afecta a la convergencia. Otras vías contemplan explorar funciones de evaluación que capturen información temporal más fina y extender la herramienta a nuevos modelos neuronales y sinápticos, aprovechando la modularidad de la arquitectura actual.
